% sample.tex with paragraphs inserted before the introduction, so all later
% text moves down: tests keeping the reading position by text on reload.
% Built by scripts/build-fixtures.sh.
\documentclass[11pt]{article}
\usepackage[a4paper,margin=25mm]{geometry}
\usepackage{amsmath,amssymb}
\usepackage{hyperref}
\title{temzdf sample document}
\author{Fixture}
\date{}
\begin{document}
\maketitle
\tableofcontents
Inserted paragraph 0: text added above the introduction to shift everything below it.\par
Inserted paragraph 1: text added above the introduction to shift everything below it.\par
Inserted paragraph 2: text added above the introduction to shift everything below it.\par
Inserted paragraph 3: text added above the introduction to shift everything below it.\par
Inserted paragraph 4: text added above the introduction to shift everything below it.\par
Inserted paragraph 5: text added above the introduction to shift everything below it.\par
\section{Introduction}\label{sec:intro}
This document exercises the viewer. The key result is Equation~\eqref{eq:gauss}
in Section~\ref{sec:math}; see also \url{https://example.com/}.
The word needle appears on several pages so that search has hits to find.
\section{Mathematics}\label{sec:math}
\begin{equation}\label{eq:gauss}
  \int_{-\infty}^{\infty} e^{-x^2}\,dx = \sqrt{\pi}
\end{equation}
\begin{align}
  \sum_{k=1}^{n} k &= \frac{n(n+1)}{2}, &
  \prod_{p\ \text{prime}} \frac{1}{1-p^{-s}} &= \zeta(s).
\end{align}
A matrix and a needle:
\[
  A = \begin{pmatrix} a_{11} & a_{12} \\ a_{21} & a_{22} \end{pmatrix},\qquad
  \det A = a_{11}a_{22} - a_{12}a_{21}.
\]
\newpage
\section{Filler}\label{sec:filler}
\subsection{First part}
\newcount\i \i=0
\loop\ifnum\i<12 \advance\i by 1
Paragraph \the\i: Lorem ipsum dolor sit amet, consectetur adipiscing elit,
sed do eiusmod tempor incididunt ut labore et dolore magna aliqua.
Ut enim ad minim veniam, quis nostrud exercitation ullamco laboris.\par
\repeat
\subsection{Second part}
Back to Section~\ref{sec:intro}. One more needle here.
\newpage
\section{Conclusion}\label{sec:conclusion}
The end, with a last needle.
\end{document}
